
\subsubsection{1.1 The Long-Context Challenge}

Transformers achieve strong performance on natural language tasks but scale quadratically with sequence length, making long-context processing computationally expensive. State Space Models (SSMs) such as Mamba offer linear-time alternatives. Rather than training SSMs from scratch, distillation transfers knowledge from pretrained Transformers with substantially reduced data requirements. MOHAWK demonstrates that Phi-Mamba can match 85\% of Phi-1.5's performance using 3B tokens.

\subsubsection{1.2 The Research Question}

A question that has not been systematically studied is which distillation objective produces representations that remain stable when sequence lengths exceed training lengths. MOHAWK employs matrix-level supervision, minimizing the discrepancy between teacher and student attention mixer outputs. CAB employs token-level supervision, aligning query/key projections (Q, K) to SSM parameters (B, C) through learned bridges.

This work tests whether token-level objectives produce representations with lower drift across sequence lengths, hypothesizing that token-level supervision is position-agnostic while matrix-level supervision encodes positional relationships that change with sequence length.

\subsubsection{1.3 Contributions}

This work makes three contributions:

First, a unified Phi-Mamba framework was constructed implementing both MOHAWK and CAB objectives in a single codebase, enabling controlled comparison.

Second, measurements show that token-level distillation produces drift slope $5\times$ lower than matrix-level distillation across sequence lengths 512--2048 tokens, with bounded drift ratio (1.34 versus 2.02).

Third, the experimental scope conditions were established: Phi-1.5 has a hard 2048-token context limit due to fixed positional embeddings, preventing analysis at longer sequences without a different teacher model.

